%% LyX 2.0.0 created this file.  For more info, see http://www.lyx.org/.
%% Do not edit unless you really know what you are doing.
\documentclass[english]{scrreprt}
\renewcommand{\familydefault}{\sfdefault}
\usepackage[T1]{fontenc}
\usepackage[latin9]{inputenc}
\setcounter{secnumdepth}{3}
\setcounter{tocdepth}{3}
\usepackage{textcomp}
\usepackage{amstext}

\makeatletter

%%%%%%%%%%%%%%%%%%%%%%%%%%%%%% LyX specific LaTeX commands.
%% Because html converters don't know tabularnewline
\providecommand{\tabularnewline}{\\}

\makeatother

\usepackage{babel}
\begin{document}

\title{ISMAR Retrieval Database}


\author{Manfred Brath}

\maketitle

\addsec{Database description}

The database consists of $6$ grids and $16$ state variables. By
the matlab function {}``ismar\_create\_database'' all the different
states will be regridded to common grids of pressure coordinates and
looking direction, so that they all will be on the same grids.


\subsubsection*{Grids }
\begin{itemize}
\item $case\_index$, an index that names a specific case with a number.
\item $pressure$, the vertical grid. In the matlab function {}``ismar\_create\_database'',
it is as first try defined from $1000\,\textrm{hPa}$ to $30\,\textrm{hPa}$
with $75$ logarithmic distributed steps. %
\footnote{In the discussion during the arts week in September 2014, we originally
decided to have also an additional flight altitude grid in pressure
units. I omitted the flight altitude grid here, because from my point
of view we do not need it. Please correct me, if I am wrong. %
}
\item $channel\_no$, the index of the channel, see table \ref{tab:Channel-definition}
for the definition.
\item $looking\_direction$, the angular grid. In the matlab function {}``ismar\_create\_database'',
it is as first try defined from $0\text{\textdegree}$ to $180\text{\textdegree}$
in steps of $5\text{\textdegree}$ ($37$ steps)
\item $scattering\_element\_index$, an index that names a specific scattering\_element
with a number.
\item $polarization$, the polarization grid (vertical, horizontal)
\end{itemize}

\subsubsection*{State variables}
\begin{itemize}
\item the channel definition, which consists of the center frequencies and
the offsets of the sidebands, see table \ref{tab:Channel-definition}.
\[
channel\_def\left(channel\_no\right)
\]
\begin{table}
\begin{centering}
\begin{tabular}{c|r@{\extracolsep{0pt}.}l|r@{\extracolsep{0pt}.}l|c}
 & \multicolumn{2}{c|}{Center fre- } & \multicolumn{2}{c|}{Side} & \tabularnewline
$channel\_no$ & \multicolumn{2}{c|}{quency {[}$\text{GHz}${]}} & \multicolumn{2}{c|}{bands {[}$\text{GHz}${]}} & Instrument\tabularnewline
\hline 
1 & \multicolumn{2}{c|}{23.8} & \multicolumn{2}{c|}{$\ \pm$0.07} & Deimos\tabularnewline
\hline 
2 & \multicolumn{2}{c|}{50.1} & \multicolumn{2}{c|}{$\pm$0.08} & Deimos\tabularnewline
\hline 
3 & \multicolumn{2}{c|}{89.0} & \multicolumn{2}{c|}{$\pm$1.1} & MARSS\tabularnewline
\hline 
4 & \multicolumn{2}{c|}{118.75} & \multicolumn{2}{c|}{$\pm$1.1} & ISMAR\tabularnewline
\hline 
5 & \multicolumn{2}{c|}{$\ $118.75} & \multicolumn{2}{c|}{$\pm$1.5} & ISMAR\tabularnewline
\hline 
6 & \multicolumn{2}{c|}{118.75} & \multicolumn{2}{c|}{$\pm$2.1} & ISMAR\tabularnewline
\hline 
7 & \multicolumn{2}{c|}{118.75} & \multicolumn{2}{c|}{$\pm$3.0} & ISMAR\tabularnewline
\hline 
8 & \multicolumn{2}{c|}{118.75} & \multicolumn{2}{c|}{$\pm$5.0} & ISMAR\tabularnewline
\hline 
9 & \multicolumn{2}{c|}{157.05} & \multicolumn{2}{c|}{$\pm$2.6} & MARSS\tabularnewline
\hline 
10 & \multicolumn{2}{c|}{183.31} & \multicolumn{2}{c|}{$\pm$1.0} & MARSS\tabularnewline
\hline 
11 & \multicolumn{2}{c|}{183.31} & \multicolumn{2}{c|}{$\pm$3.0} & MARSS\tabularnewline
\hline 
12 & \multicolumn{2}{c|}{183.31} & \multicolumn{2}{c|}{$\pm$7.0} & MARSS\tabularnewline
\hline 
13 & \multicolumn{2}{c|}{243.20} & \multicolumn{2}{c|}{$\pm$2.5} & ISMAR\tabularnewline
\hline 
14 & \multicolumn{2}{c|}{325.15} & \multicolumn{2}{c|}{$\pm$1.5} & ISMAR\tabularnewline
\hline 
15 & \multicolumn{2}{c|}{325.15} & \multicolumn{2}{c|}{$\pm$3.5} & ISMAR\tabularnewline
\hline 
16 & \multicolumn{2}{c|}{325.15} & \multicolumn{2}{c|}{$\pm$9.5} & ISMAR\tabularnewline
\hline 
17 & \multicolumn{2}{c|}{424.70} & \multicolumn{2}{c|}{$\pm$1.0} & ISMAR\tabularnewline
\hline 
18 & \multicolumn{2}{c|}{424.70} & \multicolumn{2}{c|}{$\pm$1.5} & ISMAR\tabularnewline
\hline 
19 & \multicolumn{2}{c|}{424.70} & \multicolumn{2}{c|}{$\pm$4.0} & ISMAR\tabularnewline
\hline 
20 & \multicolumn{2}{c|}{448.0} & \multicolumn{2}{c|}{$\pm$1.4} & ISMAR\tabularnewline
\hline 
21 & \multicolumn{2}{c|}{448.0} & \multicolumn{2}{c|}{$\pm$3.0} & ISMAR\tabularnewline
\hline 
22 & \multicolumn{2}{c|}{448.0} & \multicolumn{2}{c|}{$\pm$7.2} & ISMAR\tabularnewline
\hline 
23 & \multicolumn{2}{c|}{664.0} & \multicolumn{2}{c|}{$\pm$4.2} & ISMAR\tabularnewline
\hline 
24 & \multicolumn{2}{c|}{874.4} & \multicolumn{2}{c|}{$\pm$6.0} & ISMAR\tabularnewline
\hline 
\end{tabular}
\par\end{centering}

\caption{\label{tab:Channel-definition}Channel definition}
\end{table}

\item the properties of the scattering elements: 

\begin{itemize}
\item mass $mass(scattering\_element\_index)$ 
\item maximum dimension $dmax(scattering\_element\_index)$
\item volume $volume(scattering\_element\_index)$ 
\item type of scattering scattering elemenent $scattering\_element\_type(scattering\_element\_index)$,
an cell array with a tag according to the arts method \textquotedbl{}ParticleSpeciesSet\textquotedbl{}
\end{itemize}
\item the atmospheric data:

\begin{itemize}
\item temperature $temperature(case\_index,\, pressure)$ 
\item altitude $atltitude(case\_index,\, pressure)$ 
\item particle number density $pnd\_field(case\_index,\, pressure,\, scattering\_element\_index)$ 
\item brightness temperature $t\_b(case\_index,\, pressure,\, channel\_no,\, looking\_direction,\, polarization)$
\end{itemize}
\item the surface properties 

\begin{itemize}
\item surface altitude $surface\_altitude(case\_index)$
\item surface temperature $surface\_temperature(case\_index)$ 
\item surface type $surface\_type(case\_index)$, a cell array of strings
\item surface reflectivity $surface\_reflectivities(case\_index,\, channel\_no,\, looking\_direction,\, polarization,\, polarization)$,
the two last dimensions are the $2\times2$ reflection matrices of
a two dimensional stokes vector with the first component being the
vertical and the second component being the horizontal polarization.
\end{itemize}
\item the simulated position in space and time 

\begin{itemize}
\item latittude $lat(case\_index)$ 
\item longitude $lon(case\_index)$
\item time $time(case\_index)$
\end{itemize}
\end{itemize}

\addsec{Create database}

To create the database, you can use the matlab function {}``ismar\_create\_database''
within {}``atmlab/retrieval/ismar/hamburg/''. The database will
be stored as mat-files. Of course, each single case (preprocessed
ARTS-simulation) has to be a mat-file. The single preprocessed ARTS-simulations
will be combined to one database and regridded to common grids of
pressure coordinates and looking direction. Additionally, $doit\_i\_field1D\_spectrum$
will be converted to brightness temperature (Planck Tb) and then mapped
to the ISMAR channels, see also table \ref{tab:Channel-definition}.
All simulations for the database must consist of the same scattering
elements. This means all simulations to be put into the database must
have the same scattering meta data and single scattering data. Each
ARTS-simulation must consist of the into matlab converted workspace
variables:
\begin{itemize}
\item $atm\_fields\_compact$ 
\item $cloudbox\_limits$ 
\item $f\_grid$ 
\item $doit\_i\_field1D\_spectrum$ 
\item $part\_species$ 
\item $pnd\_field$ 
\item $scat\_data\_per\_part\_species$
\item $scat\_meta\_array$
\item $scat\_za\_grid$
\end{itemize}
and three additional created variables:
\begin{itemize}
\item $lat\_lon\_time$, a vector consisting of the geographical position
of the simulation (1st component is latitude, 2nd is longitude) and
the simulated time (third component) in a numeric format 
\item $surface\_properties$, an array of the two arts workspace variable
$z\_surface$ and $t\_surface$. 1st component is $z\_surface$ and
the 2nd is $t\_surface$
\item $surface\_radiation\_parameter$, which is an empty array if the surface
is a blackbody, or it is a structure of an gridded\_field6, if the
reflectivities are given, or it is a structure of a gridded\_field3
if the complex refraction indices are given. See for additional information
the ARTS documentation. \end{itemize}

\end{document}
